
A pesquisa bibliográfica foi conduzida de forma sistemática, estruturando-se em três temas principais: Hardware, Software e Redes. Para a seleção dos trabalhos, utilizaram-se as bases de dados Google Scholar, MDPI e Intechopen.

O processo de seleção seguiu um fluxo de filtragem (funil). Inicialmente, foram identificados 120 artigos por meio da análise de títulos, utilizando strings de busca majoritariamente em inglês, tais como ``iot AND energy efficiency'' e ``ESP32 AND camera monitoring''. Como ferramentas auxiliares na descoberta de conexões entre autores e na triagem inicial de relevância, foram utilizadas as plataformas ResearchRabbit e Gemini Deep Research.

Na etapa de refinamento, aplicou-se a leitura dos resumos (abstracts) e a verificação de disponibilidade integral dos textos. Foram adotados como critérios de inclusão: artigos publicados entre 2020 e 2026, relevância direta com sistemas de telemetria e foco em implementações práticas. Como critério de exclusão, descartaram-se trabalhos com acesso restrito ou que não apresentassem detalhes técnicos sobre a integração dos eixos propostos.

\begin{table}[H]
\centering
\caption{Strings de busca utilizadas no levantamento bibliográfico.}
\label{tab:strings_busca}
\begin{tabular}{@{}lll@{}}
\toprule
\textbf{ID} & \textbf{Eixo Temático} & \textbf{String de Busca} \\ \midrule
ST01 & Redes & "IoT" AND "\gls{lorawan}" AND "water waste" \\
ST02 & Contexto & smart water meter market \\
ST03 & Contexto & smart water meter costs and benefits \\
ST04 & Hardware & "IoT" AND "energy efficiency" \\
ST05 & Hardware & "ESP32" AND "camera monitoring" \\
ST06 & Hardware & "\gls{esp32cam}" AND "analog meter" AND "OCR" \\
ST08 & Software & "automatic meter reading" OCR accuracy "error rate" evaluation \\
ST09 & Redes & "packet delivery ratio" \gls{lorawan} "latency" monitoring system \\
ST010 & Contexto & "water consumption prediction" "mean absolute error" forecasting metrics \\ \bottomrule
\end{tabular}
\end{table}
